\begin{lemma}[Saturation and selected middle edges]
Let $U,W$ be finite sets of equal size $n$, let $k\in\mathbb N$, and let
$f$ be a feasible integral net flow on the network of
\noderef{BipartiteFactorNetwork} with value $kn$.
Every source-to-left and right-to-sink net flow equals $k$.
For every copied $u\in U,w\in W$, the middle net flow is an integer
between zero and one; if it is one then $uw$ is a $G$ edge.
For each copied $u$, exactly $k$ distinct copied vertices $w$ have middle
flow one, and for each copied $w$, exactly $k$ distinct copied vertices
$u$ have middle flow one.
\end{lemma}
\begin{proof}
Use the capacity table of \noderef{BipartiteFactorNetwork} and the skew
symmetry and bounds of \noderef{FiniteIntegralFlow}. If both directional
capacities are zero, both net flows are zero. If only the forward
capacity can be positive, its net flow is nonnegative, since its negative
is bounded by the zero reverse capacity. Diagonal net flows are zero by
skew symmetry. Thus the source divergence is the sum of the $n$
source-to-left flows, all in $[0,k]$. Its value is $kn$, so all these
flows equal $k$: if one were less than $k$, its positive integer deficit
would make the sum of the nonnegative deficits strictly positive,
contrary to $nk-\sum_u f(s,u)=0$.

The total divergence over all vertices is zero, since ordered-pair terms
cancel by skew symmetry. Conservation at every nonterminal therefore
says that the sink divergence is $-kn$. Only right-to-sink arcs can
contribute to its negative, so the sum of their nonnegative flows is
$kn$. There are $|W|=n$ of them, each bounded by $k$, and the same
deficit argument saturates every one.

Every middle flow is nonnegative and bounded by its capacity, which is
zero or one. Integrality makes its value zero or one. A value of one
forces capacity one, hence $G$-adjacency. For a fixed left vertex $u$,
the only potentially nonzero outgoing net flows are $f(u,s)=-k$ and
$f(u,w)$ for right vertices $w$. Its conservation equation gives
$\sum_w f(u,w)=k$. A sum of zero-or-one integers counts exactly the
distinct indices where the summand is one, giving the asserted row
count. At a right vertex $w$, conservation similarly gives
$k-\sum_u f(u,w)=0$, proving the column count. The exact subtype copies
ensure that each index represents one ambient vertex; no selected edge
is repeated. For empty parts the pointwise claims are vacuous, and for
$k=0$ the same nonnegative-sum argument gives no selected middle edges.
\end{proof}
